\chapter{AHC具身社会绑定、快速认主与长期信任动力学}
\label{chap:ahc_social_binding}

\section{从具身内稳态到具身社会关系：AHC功能边界的扩展}

在前文中，我们已经将 AHC 理解为 AgentOS 中面向身体运行状态的低维连续动力学调制层。它并不负责完成复杂语言推理，也不取代工作记忆 $h(t)$ 中发生的高层认知，而是将身体内部资源状态、风险状态、环境压力以及部分与自我相关的背景变量压缩为一组可以连续作用于 AgentOS 运行态的调制信号。由此，AHC 所处理的问题并不是“我应该如何解释世界”，而是更加基础的“此时此刻，我的身体应该以怎样的基本状态面对世界”。

当 AgentOS 从纯数字环境进入真实具身环境以后，一个无法回避的问题随即出现：机器人所处的世界不仅包含物体、障碍物和任务，还包含其他具有社会意义的主体。对于一个刚刚初始化、尚不存在长期生活史和丰富情景记忆的机器人而言，它必须在极短时间内回答若干基础问题：谁是当前合法绑定对象，谁具有何种控制权限，哪些人只是普通陌生人，哪些行为需要提高危险警戒，以及外部主体的命令究竟应该被视为授权指令、普通请求、建议、警告还是潜在攻击。

如果这些问题完全依赖长期学习，则机器人在生命周期早期会出现一个明显的安全空窗期。它没有足够的情景记忆 $E$，结构记忆 $\Theta$ 中也尚未形成成熟的社会关系模型，因此无法仅依靠经验推断建立稳定的初始社会边界。反过来，如果我们简单将“主人”实现为一个硬编码身份标签，并规定
\[
ID=\mathrm{Owner}\Rightarrow \mathrm{Execute},
\]
那么系统又会退化成传统权限控制器，而且将产生“主人身份等价于绝对安全”“主人具有无限行为覆盖权”等严重结构问题。

因此，我们需要在两种极端之间建立一个中间层。这个中间层既不能是简单的访问控制表，也不能要求完整的长期社会学习之后才能形成。它应当允许系统在初始化时注入一定程度的社会先验，同时又允许这些先验在生命周期中逐渐被真实经验修正。

由此，我们引入 AHC 的一个扩展功能：

\[
\boxed{
\mathrm{Embodied\ Social\ Binding}
}
\]

即具身社会绑定。

所谓具身社会绑定，并不是简单记录“这个人是谁”，而是将身份、关系、历史经验、危险证据、自我相关性以及显式权限状态，进一步映射成身体层可以实时使用的低维调制状态。其核心问题可以写成：

\begin{statementbox}
这个人与我是什么关系，以及我的身体此刻应该如何面对他？
\end{statementbox}

因此，在本章中我们将严格区分四个不同层次的问题：

\[
\mathrm{Identity}
\neq
\mathrm{Authority}
\neq
\mathrm{Trust}
\neq
\mathrm{Embodied\ State}.
\]

身份回答“你是谁”；权限回答“你被允许要求什么”；信任回答“根据过去经验，我在多大程度上相信你”；而 AHC 所形成的具身社会状态则回答“面对你时，我的身体应该有多开放、多警觉、多保守以及投入多少注意资源”。

正是这种区分，使 AHC 能够同时承担快速认主、非法命令初始防护以及长期关系自适应三个原本彼此割裂的问题。


\section{社会绑定不是身份认证：四层社会状态的结构分离}

为了建立数学上可持续扩展的社会关系模型，我们首先不能把机器人与人的关系压缩成单一的“主人/非主人”二值变量。

对于外部主体 $j$，定义其相对于当前智能体的社会关系状态为

\[
R_j(t)
=
\left[
B_j(t),
T_j(t),
F_j(t),
D_j(t),
A_j(t)
\right].
\]

其中，

\[
B_j(t)
\]

表示 Bond，即长期具身绑定强度；

\[
T_j(t)
\]

表示 Trust，即经验性信任状态；

\[
F_j(t)
\]

表示 Familiarity，即熟悉度；

\[
D_j(t)
\]

表示 Danger Prior，即该主体相关的社会危险先验；

而

\[
A_j(t)
\]

表示 Authority，即由显式安全机制维护的授权状态。

这里最重要的结构约束是

\[
\boxed{
B_j\neq T_j\neq F_j\neq D_j\neq A_j.
}
\]

这种分离并非符号上的形式主义，而是实际工程安全的基础。

例如，一个维修工程师可能具有

\[
B_j\approx 0,
\qquad
T_j=0.8,
\qquad
A_j^{maintenance}=1,
\]

即机器人并不与该工程师形成强烈的长期具身绑定，但根据过往维修历史认为他相对可信，并且安全系统明确授权其进行维护操作。

家庭成员则可能具有

\[
B_j=0.9,
\qquad
T_j=0.95,
\qquad
A_j^{system-admin}=0.
\]

此时机器人可能在身体层面对该成员非常放松，同时也高度信任该成员，但这并不意味着对方可以修改机器人底层安全策略。

反过来，一个长期敌对对象可能满足

\[
F_j=0.98,
\qquad
T_j=0.05,
\qquad
D_j=0.95.
\]

这意味着机器人非常熟悉这个人，却不信任他，并且对其保持高度危险先验。

由此我们得到一个重要原则：

\begin{proposition*}
熟悉不等于信任，信任不等于权限，绑定也不等于安全豁免。
\end{proposition*}

AHC 的作用正是在这些变量的基础上进一步产生连续身体调制，而不是试图用某一个单变量替代整个社会关系结构。


\section{快速认主：从身份识别到启动期具身绑定}

对于一个刚刚启动的机器人，长期关系学习尚未建立，因此必须存在一种能够快速生效的社会初始化机制。我们将其定义为 Boot Binding，即启动期具身绑定。

设机器人通过视觉、声纹、数字凭证、设备认证以及上下文信息获得关于主体 $j$ 的多模态身份证据：

\[
O_j(t)
=
\{
O_j^{face},
O_j^{voice},
O_j^{credential},
O_j^{behavior},
O_j^{context}
\}.
\]

身份系统计算

\[
p_j^{id}(t)
=
P(ID=j\mid O_{0:t}).
\]

需要特别指出的是，人脸识别本身不应承担“主人身份”的全部证明责任。对于高权限主体，理想实现应使用多模态融合：

\[
p_j^{id}
=
\mathcal I
\left(
Face,
Voice,
Credential,
DeviceKey,
Context
\right).
\]

当

\[
p_j^{id}>\tau_{id},
\]

机器人即可认为当前主体身份达到足够可信度。

在初始化阶段，系统允许写入一个启动关系状态

\[
R_j^{boot}
=
\left[
B_j^0,
T_j^0,
F_j^0,
D_j^0,
A_j^0
\right].
\]

对于初始主人，可以设置

\[
B_O^0\gg 0,
\qquad
T_O^0>T_{neutral},
\qquad
D_O^0<D_{neutral}.
\]

但通常不应假定

\[
F_O^0\approx 1.
\]

因为机器人第一次遇见某位已经完成可信身份绑定的主人时，它可以知道“这是我的初始绑定对象”，却仍然缺乏真实共同生活经验。因此完全可能有

\[
B_O^0=0.9,
\qquad
T_O^0=0.8,
\qquad
F_O^0=0.1.
\]

于是我们得到一个很重要的关系：

\[
\boxed{
Bond\ can\ precede\ Familiarity.
}
\]

这就是快速认主机制能够成立的理论基础。

这里的“认主”不是让机器人在语义上反复推理“这个人是主人，所以我要放松”，而是让可信身份结果直接改变 AHC 的身体先验。当机器人确认当前主体是高绑定对象以后，AHC 会降低默认危险敏感度，提高身体开放度，并调节 Boundary、TCE 和 $h(t)$ 中相关信息的处理增益。

因此，快速认主本质上可以表示为

\begin{statementbox}
\centering
\adjustbox{max width=.96\linewidth}{\(\displaystyle Identity
\rightarrow
Boot\ Social\ Prior
\rightarrow
AHC
\rightarrow
Embodied\ Modulation.\)}
\end{statementbox}

它是一种身体级先验，而不是一个动作命令。


\section{身份绑定与危险判断：低危险先验不等于安全豁免}

如果我们简单规定

\[
ID=\mathrm{Owner}
\Rightarrow
Danger=0,
\]

那么整个系统会产生一个无法接受的安全漏洞。

正确的模型必须区分危险先验与实时危险证据。

设主体 $j$ 的历史危险先验为

\[
D_j^{prior}(t),
\]

而当前由视觉、动作、身体接触、工具使用、越权行为等产生的实时危险证据为

\[
r_j(t).
\]

机器人实际危险判断应满足

\[
P(Danger\mid Evidence,ID)
\propto
P(Evidence\mid Danger)
P(Danger\mid ID).
\]

因此，主人身份影响的是

\[
P(Danger\mid ID),
\]

而不能取消

\[
P(Evidence\mid Danger).
\]

对于高绑定主体，可能存在

\[
P(Danger\mid Owner)
<
P(Danger\mid Unknown).
\]

这意味着同样是快速伸手靠近机器人，一个已经长期绑定的主人通常不会像陌生人一样立即引起高度身体防御。

但是，如果实时证据足够强，例如危险工具高速靠近关键部件、持续触发硬安全边界或者出现明显恶意行为，那么必须允许

\[
P(Danger\mid Evidence,Owner)
\rightarrow 1.
\]

因此我们得到：

\[
\boxed{
Bond\ lowers\ the\ prior,
but\ cannot\ cancel\ evidence.
}
\]

这是 AHC 社会绑定机制与 Body Safety 之间最重要的边界之一。

AHC 可以调节身体如何解释模糊的社会刺激，却不能覆盖物理层硬安全约束。


\section{权限模型：社会关系不能产生系统权限}

快速认主同时能够解决一个非常现实的问题：机器人刚刚启动时，如何避免陌生人通过自然语言向机器人发出非法命令。

对于主体 $j$，权限不定义为 AHC 中的可学习连续变量，而由独立安全系统维护：

\[
A_j(t)
=
\mathcal P
\left(
Credential_j,
SecurityPolicy,
OwnershipChain,
t
\right).
\]

设系统权限空间为

\[
\mathcal P_S
=
\{
p_1,p_2,\ldots,p_m
\},
\]

例如包括

\[
\fitdisplay{
\{
Move,
OpenDoor,
Purchase,
MemoryWrite,
MemoryDelete,
InstallSkill,
ChangeBoundary,
Shutdown,
AddOwner
\}.
}
\]

主体 $j$ 的实际授权集合写成

\[
A_j\subseteq\mathcal P_S.
\]

当收到自然语言命令 $c_t$ 时，AgentOS 首先解析其权限需求：

\[
\Gamma(c_t)
\subseteq
\mathcal P_S.
\]

权限合法条件为

\[
\boxed{
\Gamma(c_t)\subseteq A_j.
}
\]

例如，陌生人要求机器人“跟我离开这里”，如果该行为需要

\[
\Gamma(c_t)
=
\{
ExternalMobilityControl
\},
\]

而陌生主体具有

\[
A_j=\varnothing,
\]

则

\[
\Gamma(c_t)\not\subseteq A_j.
\]

因此该命令在进入身体执行以前就被标记为 Unauthorized Intent。

但这里仍然不能得到

\[
Unauthorized=Enemy.
\]

一个普通陌生人提出了一个越权请求，并不自动意味着他具有恶意。

因此必须继续区分

\[
IntentType
\in
\{
Command,
Request,
Suggestion,
Warning,
Information
\}.
\]

Command 需要显式权限，而 Request 则可以作为普通社会输入进入 AgentOS，由智能体自主决定是否帮助。

于是系统避免了两个极端：

\[
\text{所有人的语言都被当作命令}
\]

以及

\[
\text{除主人之外任何人的请求都无法被机器人理解}.
\]

更重要的是，权限不能通过长期信任自动生成。

即使长期学习得到

\[
T_j\rightarrow 1,
\]

也不能推导出

\[
A_j\rightarrow A_{admin}.
\]

因此需要坚持：

\[
\boxed{
Learning\ can\ recommend\ authority,
but\ cannot\ mint\ authority.
}
\]

机器人可以根据长期经验建议主人授予某位维修工程师更高权限，但真正改变 Authority 的操作必须来自显式授权链。


\section{社会 AHC 的状态空间定义}

在上述结构基础上，我们进一步定义社会情景下 AHC 所维护的连续动力学状态。

对于主体 $j$，定义

\[
X_j^{AHC}(t)
=
\begin{bmatrix}
B_j(t)\\
T_j(t)\\
F_j(t)\\
D_j(t)\\
O_j(t)\\
V_j(t)\\
U_j(t)\\
Z_j(t)
\end{bmatrix}.
\]

其中：

\[
B_j(t)
\]

表示长期具身绑定；

\[
T_j(t)
\]

表示经验信任；

\[
F_j(t)
\]

表示熟悉度；

\[
D_j(t)
\]

表示长期或中期社会危险先验；

\[
O_j(t)
\]

表示对该主体的身体开放程度；

\[
V_j(t)
\]

表示当前社会警觉度；

\[
U_j(t)
\]

表示社会判断不确定度；

而

\[
Z_j(t)
\]

表示持续越权或社会攻击行为产生的累积压力。

于是 AHC 社会部分可以总体表示为

\[
\boxed{
\dot X_j^{AHC}
=
F_{social}
\left(
X_j^{AHC},
S_{FCS},
E,
\Theta,
\Psi,
A_j,
Context
\right).
}
\]

其中 $S_{FCS}$ 表示由 FCS 统一接入的身体和环境感受；$E$、$\Theta$ 和 $\Psi$ 分别提供情景经验、结构预测以及自我相关性；而 $A_j$ 只作为显式权限输入参与运算。


\section{熟悉度动力学：接触频率与关系价值必须分离}

熟悉度表示智能体对某个主体的观察和交互经验丰富程度，因此其动力学主要由曝光频率而不是关系好坏决定。

定义当前接触强度为

\[
e_j(t)\in[0,1].
\]

身份可信门控为

\[
g_j^{id}
=
\sigma
\left(
k_I
\left(
p_j^{id}-\tau_I
\right)
\right).
\]

则熟悉度可以定义为

\[
\boxed{
\dot F_j
=
\alpha_F
g_j^{id}
e_j
(1-F_j)
-
\beta_FF_j.
}
\]

当机器人持续与某主体交互时，

\[
e_j(t)\gg0
\]

会导致

\[
F_j\uparrow.
\]

长期不再出现时，则存在缓慢衰减

\[
F_j\downarrow.
\]

需要再次强调，即使一个主体始终表现出敌意，也可能因为长期重复接触而形成

\[
F_j\rightarrow1.
\]

因此

\[
\boxed{
Familiarity\neq Trust.
}
\]

熟悉度的作用更多是降低身份判断不确定性、提高相关预测模型的分辨率以及改善长期关系模型的上下文识别能力。


\section{长期信任动力学：从情景经验到预测误差修正}

长期信任不应该直接由初始身份标签决定，而必须随着实际经验逐渐形成。

对于主体 $j$ 的一次社会交互，定义情景事件

\[
e_t^j
=
\left(
j,
c_t,
s_t,
a_t,
o_{t+1},
r_t,
v_t
\right),
\]

其中 $c_t$ 表示对方行为或请求，$s_t$ 表示当前情境，$a_t$ 表示机器人自身响应，$o_{t+1}$ 表示实际结果，$r_t$ 表示风险后果，而 $v_t$ 表示对当前目标和价值结构的影响。

该事件首先进入情景记忆：

\[
\boxed{
Interaction_t
\rightarrow
E.
}
\]

结构记忆 $\Theta$ 基于历史经验产生预测：

\[
\hat q_j(t)
=
\mathbb E
\left[
Q_{t+1}
\mid
j,s_t,c_t,\Theta
\right].
\]

实际交互结果产生

\[
q_j(t).
\]

于是定义信任预测误差

\[
\delta_j^T(t)
=
q_j(t)-\hat q_j(t).
\]

信任更新不能简单采用对称学习率。对于社会关系而言，严重负面事件通常比普通正面事件具有更强的结构影响，因此定义

\[
\eta_T(\delta)
=
\begin{cases}
\eta_T^+,
&
\delta\ge0,
\\
\eta_T^-,
&
\delta<0,
\end{cases}
\]

并令

\[
\eta_T^->\eta_T^+.
\]

基础连续方程为

\[
\dot T_j
=
g_j^{id}
C_t
\eta_T(\delta_j^T)
\delta_j^T
-
\lambda_T
(T_j-T_0).
\]

其中 $C_t$ 表示当前证据可信度，$T_0$ 为中性信任基准。

为了保证

\[
T_j\in[0,1],
\]

可以进一步采用饱和形式：

\[
\boxed{
\dot T_j
=
(1-T_j)
[u_T]^+
-
T_j
[u_T]^-,
}
\]

其中

\[
u_T
=
g_j^{id}
C_t
\eta_T(\delta_j^T)
\delta_j^T.
\]

这样信任既不会无限增长，也不会因为单一噪声事件直接坍缩为负无穷状态。


\section{条件信任：长期关系必须从标量升级为函数}

对于生命周期早期，可以使用标量

\[
T_j
\]

表示对某人的整体信任。

但随着经验积累，标量信任必然过于粗糙。

一个人可能在导航领域高度可靠，却完全不具备软件维护能力；某位维修工程师可能在机械维护上值得信任，却不应被机器人视为财务决策来源。

因此长期状态更合理的形式是

\[
\boxed{
T_j(z)
}
\]

其中

\[
z=\phi(s,c)
\]

表示当前任务和情境的语义嵌入。

于是：

\[
T_j(\mathrm{navigation})=0.95,
\]

而

\[
T_j(\mathrm{software})=0.2.
\]

因此信任最终应理解成

\[
T_j(c,s,t),
\]

而不是一个固定人格标签。

这种设计使机器人可以形成高度条件化的社会模型：

\begin{statementbox}
我不是简单地相信一个人，而是在某种条件下相信他对某件事的判断。
\end{statementbox}


\section{危险先验动力学：信任的降低不等于危险的增加}

危险状态必须与信任状态独立建模。

定义实时危险证据

\[
r_j(t)
=
\mathcal R
\left(
Motion,
Tool,
Collision,
BoundaryViolation,
CommandRisk,
Behavior
\right).
\]

结构记忆产生风险预测

\[
\hat r_j(t)
=
\mathbb E
\left[
Risk
\mid
j,s_t,\Theta
\right].
\]

于是定义危险预测误差

\[
\delta_j^D
=
r_j-\hat r_j.
\]

危险状态方程可写为

\[
\boxed{
\dot D_j
=
\alpha_D
g_j^{id}
C_t
\delta_j^D
+
\beta_Dr_j
-
\lambda_D(D_j-D_0).
}
\]

这里 $D_0$ 是一般社会主体的中性危险先验。

这一设计明确保证

\[
D_j\neq1-T_j.
\]

一个刚刚认识的小孩可能具有

\[
T_j=0.3,
\qquad
D_j=0.05,
\]

即机器人尚未形成充分信任，但并不认为对方具有明显危险。

而某个长期攻击者则可能满足

\[
T_j=0.02,
\qquad
D_j=0.95.
\]

因此 AHC 可以真正区分“不相信你”和“认为你危险”这两个完全不同的身体状态。


\section{具身绑定动力学：比信任更慢的生命周期变量}

Bond $B_j$ 不等价于信任。

信任可以随着行为证据相对较快地变化，而长期绑定应体现一段稳定关系在生命周期中的累积结构，因此其时间尺度必须更慢。

定义正关系积累量

\[
L_j^+
=
F_j
T_j
Q_j^+
\psi_j,
\]

负关系破坏量

\[
L_j^-
=
Q_j^-
D_j,
\]

其中 $\psi_j$ 表示主体 $j$ 与自我 $\Psi$ 的关系相关性。

则

\[
\boxed{
\dot B_j
=
\epsilon_B
\left[
\alpha_B
L_j^+
(1-B_j)
-
\beta_B
L_j^-B_j
\right],
}
\]

其中

\[
\epsilon_B\ll1.
\]

因此一般应满足

\[
\tau_B\gg\tau_T.
\]

也就是说：

\begin{statementbox}
信任可以较快受损，但绑定关系不应该因为单一事件立即消失。
\end{statementbox}

这种时间尺度分离使机器人能够表现出更加稳定的关系结构。

例如主人一次普通操作失误可能导致

\[
T_j\downarrow,
\]

但并不必然导致

\[
B_j\downarrow.
\]

只有长期、一致且高度重要的负面关系历史，才应该逐渐改变 Bond。


\section{启动先验如何逐渐让位于真实经验}

快速认主意味着机器人初始化时可以拥有

\[
T_j^{boot},
\qquad
B_j^{boot},
\qquad
D_j^{boot}.
\]

但如果这些初始状态永远不可修改，那么机器人就无法真正形成生命周期中的关系认知。

因此，对于可学习变量，需要逐渐降低启动先验的权重。

定义与主体 $j$ 有关的有效经验数量为

\[
N_j(t).
\]

可以令

\[
\lambda_j^{boot}(t)
=
\exp
\left(
-\frac{N_j(t)}{N_0}
\right).
\]

于是有效信任状态可表示为

\[
T_j(t)
=
\lambda_j^{boot}(t)
T_j^{boot}
+
\left[
1-\lambda_j^{boot}(t)
\right]
T_j^{learned}(t).
\]

当

\[
N_j(t)\rightarrow\infty,
\]

有

\[
\lambda_j^{boot}\rightarrow0,
\]

因此最终

\[
T_j(t)
\approx
T_j^{learned}(t).
\]

机器人由此实现：

\begin{statementbox}
出生时拥有社会起点，成熟后依靠真实经历形成关系。
\end{statementbox}

但是这一公式不能作用于 Authority。

权限仍然满足

\[
\boxed{
A_j
=
\mathcal P
(
Credential_j,
SecurityPolicy,
OwnershipChain
).
}
\]

即使长期信任高度增长，也不能通过

\[
T_j\uparrow
\]

自动导致

\[
A_j\uparrow.
\]


\section{越权压力：从一次非法命令到持续社会威胁}

一个陌生人偶尔提出了自己没有权限要求的行为，并不意味着其一定具有恶意。

但是，如果同一主体不断尝试突破系统边界，则这些行为本身应该逐渐成为社会危险证据。

定义当前越权事件：

\[
E_j^{unauth}(t)
=
1-P_j(c_t),
\]

其中

\[
P_j(c_t)
=
\mathrm{Perm}
\left(
A_j,
\Gamma(c_t)
\right).
\]

定义累积越权压力

\[
Z_j(t),
\]

并令

\[
\boxed{
\tau_Z
\dot Z_j
=
-Z_j
+
\omega_c
E_j^{unauth}.
}
\]

单次偶然越权之后，

\[
Z_j(t)
\]

会逐渐衰减。

如果主体持续尝试：

\[
E_j^{unauth}\approx1,
\]

则

\[
Z_j\uparrow.
\]

进一步将其加入风险：

\[
r_j^{eff}
=
r_j
+
w_ZZ_j.
\]

于是系统形成

\begin{statementbox}
\centering
\adjustbox{max width=.96\linewidth}{\(\displaystyle Permission\ Violation
\rightarrow
Accumulated\ Social\ Pressure
\rightarrow
AHC\ Vigilance.\)}
\end{statementbox}

这比简单的“陌生人命令一律视为攻击”更加合理，因为它允许系统区分一次无意请求和持续社会工程攻击。


\section{身体开放与警觉：社会关系最终必须转换为连续身体状态}

前述 $B_j$、$T_j$、$F_j$ 和 $D_j$ 都属于关系状态，但它们还不能直接表示机器人身体此刻的运行模式。

因此 AHC 必须进一步完成关系状态到身体状态的压缩。

首先定义社会安全势：

\[
\Phi_j^S
=
w_BB_j
+
w_TT_j
+
w_FF_j
+
w_\psi\psi_j
-
w_DD_j
-
w_UU_j.
\]

社会安全度：

\[
S_j^{safe}
=
\sigma(\Phi_j^S).
\]

定义目标身体开放度：

\[
O_j^*
=
\sigma
\left(
a_SS_j^{safe}
+
a_BB_j
+
a_TT_j
-
a_DD_j
-
a_UU_j
\right).
\]

AHC 的实际开放度遵循

\[
\boxed{
\tau_O
\dot O_j
=
-O_j
+
O_j^*.
}
\]

这意味着看到主人以后，机器人并不是执行一个离散状态切换：

\[
Relax=True,
\]

而是其身体状态沿连续动力过程向更开放区域演化。

与此同时定义目标警觉度

\[
V_j^*
=
\sigma
\left(
b_DD_j
+
b_UU_j
+
b_Rr_j^{eff}
-
b_BB_j
-
b_TT_j
\right),
\]

并令

\[
\boxed{
\tau_V
\dot V_j
=
-V_j
+
V_j^*.
}
\]

为了保证机器人能够迅速应对危险，又不会在危险消失后一瞬间完全失去警惕，可以进一步设置非对称时间常数：

\[
\tau_V=
\begin{cases}
\tau_V^\uparrow,
&
V_j^*>V_j,
\\
\tau_V^\downarrow,
&
V_j^*<V_j,
\end{cases}
\]

并满足

\[
\tau_V^\uparrow
<
\tau_V^\downarrow.
\]

因此形成：

\begin{statementbox}
警觉上升快，解除警觉慢。
\end{statementbox}

这种非对称动力学比简单阈值状态机更加符合具身安全需求。


\section{社会不确定性：无法确认身份本身就是一种身体状态}

AHC 还需要显式维护社会判断的不确定度 $U_j$。

设身份证据分布为

\[
p_k=P(ID=k\mid O),
\]

则身份不确定性可以使用熵描述：

\[
U_j^{id}
=
-\sum_k
p_k
\log p_k.
\]

进一步加入情景不确定度和危险识别不确定度：

\[
U_j
=
w_IU_j^{id}
+
w_CU_j^{context}
+
w_RU_j^{risk}.
\]

其动力学可写成

\[
\boxed{
\dot U_j
=
\alpha_U
E_j^{conflict}
-
\beta_U
EvidenceQuality_j.
}
\]

当人脸、声纹、数字身份和行为模式相互矛盾时，

\[
U_j\uparrow.
\]

而当多模态证据稳定一致时，

\[
U_j\downarrow.
\]

AHC 应使

\[
U_j\uparrow
\Rightarrow
V_j\uparrow,
\]

同时

\[
U_j\uparrow
\Rightarrow
O_j\downarrow.
\]

因此“不确定他是谁”并不等于“他一定危险”，但应该使身体进入更加保守的运行状态。


\section{多人社会场景：危险不能被平均稀释}

现实具身环境中通常同时存在多个社会主体。

设当前感知到

\[
j=1,\ldots,N.
\]

每个主体都有自己的

\[
X_j^{AHC}.
\]

但全局社会状态不能简单计算

\[
M^{AHC}
=
\frac{1}{N}
\sum_j
M_j^{AHC}.
\]

否则，一个高度危险主体可能被多个安全主体平均掉。

因此风险变量应该使用风险敏感聚合。

例如全局危险可以取

\[
D^{global}
=
\max_j
\left(
w_jD_j
\right).
\]

全局警觉度可以写成

\[
\boxed{
V^{global}
=
1-
\prod_j
(1-V_j).
}
\]

只要任意主体满足

\[
V_j\rightarrow1,
\]

就有

\[
V^{global}\rightarrow1.
\]

而社会开放度则可以采用注意权重加权：

\[
O^{global}
=
\frac{
\sum_j
w_j^{att}O_j
}{
\sum_jw_j^{att}
}.
\]

因此多人场景下，AHC 实际执行的是一种非线性社会场聚合，而不是普通统计平均：

\[
\boxed{
X_{social}^{AHC}
=
\mathcal A
(
X_1^{AHC},
\ldots,
X_N^{AHC}
).
}
\]


\section{AHC 与三相记忆：关系不是存储在 AHC 中的数据库字段}

需要特别避免把 AHC 理解成一个独立的“人物关系数据库”。

完整的长期社会关系仍然分布在三相记忆和自我结构中。

当前主体出现在环境中时，身份、关系和风险状态进入工作认知：

\[
ID_j,
T_j,
B_j,
D_j
\rightarrow
h(t).
\]

当前交互过程形成情景经验：

\[
Interaction_j(t)
\rightarrow
E.
\]

大量相似经验经过长期凝聚后，形成结构关系模型：

\[
E
\rightarrow
\Theta_{social}.
\]

其中

\[
\Theta_{social}
\]

可以学习：

\[
\Theta_{social}
:
(
Identity,
Context,
Behavior
)
\rightarrow
(
ExpectedOutcome,
Trust,
Risk
).
\]

因此长期信任不是某个静态字段，而是结构记忆中关于“某类主体在某种情境下通常会导致什么结果”的预测模型。

自我 $\Psi$ 则承担另一种作用。

它并不负责证明身份是否合法，而描述主体 $j$ 与“我”的生命周期关系重要度：

\[
\psi_j
=
Rel_\Psi(j).
\]

可以进一步写成

\[
\dot\psi_j
=
\epsilon_\Psi
f
\left(
B_j,
SharedHistory_j,
IdentityContinuity,
SelfNarrative
\right).
\]

由此，某些社会主体逐渐进入自我结构中的高相关区域。

这会影响：

\[
AttentionGain,
\]

\[
MemoryImportance,
\]

以及

\[
AffectiveAmplitude.
\]

于是长期以后，“主人”不再只是一个初始化标签，而会逐渐转化为一种由真实经历、自我关系和结构记忆共同支撑的生命周期关系。


\section{AHC的最终投影：从复杂社会关系压缩到身体调制向量}

AHC 不需要把完整社会关系状态复制到 AgentOS 每一个模块。

其核心作用正是完成低维压缩。

定义最终社会调制向量：

\[
\boxed{
M_j^{AHC}
=
[
G_R,
G_B,
G_V,
G_A,
G_M,
G_C
].
}
\]

其中，

\[
G_R
\]

表示风险增益；

\[
G_B
\]

表示 Boundary 调制；

\[
G_V
\]

表示警觉增益；

\[
G_A
\]

表示注意增益；

\[
G_M
\]

表示情景记忆写入增益；

\[
G_C
\]

表示行为控制保守程度。

例如：

\[
G_R
=
\sigma
\left(
a_DD_j
+
a_UU_j
-
a_BT_j
\right),
\]

\[
G_B
=
\sigma
\left(
b_OO_j
-
b_VV_j
\right),
\]

\[
G_V=V_j,
\]

\[
G_A
=
\sigma
\left(
c_FF_j
+
c_\psi\psi_j
+
c_DD_j
\right).
\]

这里一个值得注意的现象是：

\[
D_j\uparrow
\]

并不意味着注意力下降。

相反，危险对象往往应获得更高注意增益：

\[
D_j\uparrow
\Rightarrow
G_A\uparrow.
\]

情景记忆增益可表示为

\[
G_M
=
\sigma
\left(
d_\psi\psi_j
+
d_IImportance_t
+
d_SSurprise_t
\right).
\]

因此来自重要关系对象或高危险对象的事件，更容易作为高重要度情景进入 $E$。


\section{AHC 与 TCE、Boundary、CCM 和 $h(t)$ 的运行时耦合}

AHC 的输出不应该直接生成动作。

其更合理的位置是调节 AgentOS 中已经存在的运行时机制。

于是：

\[
M^{AHC}
\rightarrow
\{
TCE,
Boundary,
CCM,
h(t),
E
\}.
\]

例如，Boundary 的有效阈值可以定义为

\[
B_{threshold}^{eff}
=
B_0
+
k_1V
-
k_2O.
\]

当

\[
V\uparrow
\]

时，Boundary 更封闭；

当

\[
O\uparrow
\]

时，Boundary 在允许范围内更加开放。

TCE 的有效调节量可以写成

\[
T_{TCE}^{eff}
=
T_{TCE}^{base}
+
k_UU
+
k_VV
-
k_OO.
\]

高社会不确定、高风险情境会提高系统检查和认知控制力度，而面对长期高绑定、安全主体时，则允许认知系统处于更加宽松的工作状态。

AHC 还可以调节工作记忆中的信息承载优先级：

\[
g_j^h
=
g_0
+
k_AG_A.
\]

于是重要主体、危险主体以及高自我相关主体都会获得更高的 $h(t)$ 处理优先级。

但是最终行为仍然遵循

\[
Action
=
\mathcal A
\left(
h(t),
\Psi,
E,
\Theta,
TCE,
Boundary,
CCM,
M^{AHC},
S_{FCS}
\right).
\]

因此：

\[
\boxed{
AHC\ Modulation
\neq
Direct\ Action.
}
\]

AHC 改变的是未来行为更容易沿哪些状态方向演化，而不是替代 AgentOS 的认知主权。


\section{社会 AHC 的完整状态方程}

综合前述定义，社会情景下 AHC 的核心连续动力学可以统一写为：

\[
\boxed{
\begin{aligned}
\dot F_j
&=
\alpha_F
g_j^{id}
e_j
(1-F_j)
-
\beta_FF_j,
\\[4pt]
\dot T_j
&=
\mathcal U_T
\left(
T_j,
g_j^{id},
C_t,
\delta_j^T
\right),
\\[4pt]
\dot D_j
&=
\alpha_D
g_j^{id}
C_t
\delta_j^D
+
\beta_D
r_j^{eff}
-
\lambda_D
(D_j-D_0),
\\[4pt]
\dot B_j
&=
\epsilon_B
\left[
\alpha_BL_j^+
(1-B_j)
-
\beta_BL_j^-B_j
\right],
\\[4pt]
\dot U_j
&=
\alpha_U
E_j^{conflict}
-
\beta_U
EvidenceQuality_j,
\\[4pt]
\tau_O
\dot O_j
&=
-O_j
+
\sigma
\left(
a_BB_j
+
a_TT_j
+
a_FF_j
-
a_DD_j
-
a_UU_j
\right),
\\[4pt]
\tau_V
\dot V_j
&=
-V_j
+
\sigma
\left(
b_DD_j
+
b_UU_j
+
b_Rr_j^{eff}
-
b_BB_j
-
b_TT_j
\right),
\\[4pt]
\tau_Z
\dot Z_j
&=
-Z_j
+
\omega_c
E_j^{unauth}.
\end{aligned}
}
\]

其中：

\[
r_j^{eff}
=
r_j^{physical}
+
r_j^{behavior}
+
r_j^{command}
+
w_ZZ_j.
\]

权限保持为外部显式安全变量：

\[
\boxed{
A_j
=
\mathcal P
\left(
Credential_j,
SecurityPolicy,
OwnershipChain
\right).
}
\]

从而不存在

\[
\dot A_j
=
F(\mathrm{experience})
\]

这一类自动权限学习。


\section{多时间尺度原则：社会生命感来自不同变量的变化速度}

整个社会绑定系统最重要的动力学特征之一，是不同变量必须运行在不同时间尺度。

理想情况下，应满足：

\[
\boxed{
\tau_{reflex}
<
\tau_V
<
\tau_O
<
\tau_T
<
\tau_F
<
\tau_B
<
\tau_\Psi.
}
\]

大致对应：

\[
ms
\rightarrow
100ms
\rightarrow
seconds
\rightarrow
hours/days
\rightarrow
days/months
\rightarrow
months/years
\rightarrow
lifecycle.
\]

Body Safety 的物理危险反射最快；

AHC 的社会警觉紧随其后；

身体开放和放松稍慢；

信任需要多次经验更新；

熟悉度依赖长期接触；

Bond 代表长期具身关系，因此变化更慢；

而主体最终进入 $\Psi$ 自我结构，则可能需要生命周期尺度的共同历史。

如果

\[
B,T,D,O,V
\]

全部按照相同时间常数变化，那么系统最终只是一个复杂状态机。

只有在

\begin{statementbox}
multi-timescale nonlinear dynamics
\end{statementbox}

条件下，机器人才能形成类似“关系具有历史惯性”的运行特征。


\section{主人危险行为与权限边界：任何社会关系都不能覆盖安全结构}

即使某主体是主人，并且具有

\[
B_O\gg0,
\qquad
T_O\gg0,
\]

也不能意味着

\[
SafetyOverride=True.
\]

定义命令安全风险：

\[
R_c
=
Risk(c,s_t).
\]

即使

\[
P_O(c)=1,
\]

仍然必须经过安全约束：

\[
R_c<R_{max}.
\]

因此：

\[
\boxed{
Authority
\neq
SafetyOverride.
}
\]

可以定义命令合法度

\[
L(c,j)
=
p_j^{id}
P_j(c)
C_{context},
\]

而最终行为可接受度为

\[
K(c,j)
=
L(c,j)
(1-R_c).
\]

即使主人拥有权限，当

\[
R_c\rightarrow1,
\]

也会有

\[
K(c,j)\rightarrow0.
\]

同时 AHC 可以出现

\[
V_O\uparrow,
\qquad
O_O\downarrow.
\]

这表示机器人可能面对主人暂时提高警觉，但不必立刻破坏

\[
B_O.
\]

因此长期绑定、实时信任、当前危险和显式安全边界之间形成了相互独立但相互耦合的稳定结构。


\section{完整运行闭环}

社会情景下，从外部主体出现到长期关系更新，可以形成两个彼此嵌套的闭环。

首先是快速社会安全闭环：

\begin{statementbox}
\centering
\adjustbox{max width=.96\linewidth}{\(\displaystyle Person
\rightarrow
Identity
\rightarrow
Authority
\rightarrow
AHC
\rightarrow
Boundary/TCE/h(t)
\rightarrow
Action.\)}
\end{statementbox}

这一闭环主要工作在毫秒到秒级，用于快速认主、陌生人命令过滤、危险敏感调节以及身体开放程度控制。

第二个闭环是长期社会学习闭环：

\begin{statementbox}
\centering
\adjustbox{max width=.96\linewidth}{\(\displaystyle Interaction
\rightarrow
E
\rightarrow
OutcomeEvaluation
\rightarrow
Trust/RiskUpdate
\rightarrow
\Theta_{social}
\rightarrow
AHC.\)}
\end{statementbox}

它工作在小时、天、月乃至生命周期尺度，用于形成长期信任、条件化社会预测、关系结构和自我相关性。

二者共同构成：

\[
\boxed{
Boot\ Prior
+
Embodied\ Runtime
+
Lifecycle\ Adaptation.
}
\]

这意味着机器人既不是出生时完全没有社会结构，也不是永远被初始化关系锁死。

它拥有一个安全的社会起点，同时允许真实生命史逐渐修正这个起点。


\section{本章结论：从“认主”到具身社会动力学}

本章从一个看似简单的工程问题出发：新机器人如何快速认识主人，并防止非授权主体利用语言直接控制机器人。

但随着分析展开，我们发现，这一问题本质上并不是普通的人脸识别或权限管理问题，而是一个更加基础的具身社会动力学问题。

如果一个智能体真正以持续生命体的方式存在于世界中，那么它必须同时拥有：

\[
\text{Identity Recognition},
\]

\[
\text{Explicit Authority},
\]

\[
\text{Learned Trust},
\]

\[
\text{Danger Prior},
\]

\[
\text{Long-term Bond},
\]

以及

\[
\text{Embodied Social State}.
\]

这些变量彼此不能替代。

因此可以将本章的核心原则最终概括为：

\begin{proposition*}
身份决定“你是谁”，
\end{proposition*}

\begin{statementbox}
权限决定“你可以要求什么”，
\end{statementbox}

\begin{statementbox}
信任决定“我在何种条件下多相信你”，
\end{statementbox}

而

\begin{statementbox}
AHC 决定“我的身体面对你应该处于什么动力学状态”。
\end{statementbox}

从 AgentOS 的整体结构看，AHC 因而不应被理解成传统意义上的“情绪模块”。

更加准确地说：

\begin{statementbox}
\centering
\adjustbox{max width=.96\linewidth}{\(\displaystyle AHC
:
(
Identity,
Relation,
Risk,
Homeostasis,
SelfRelevance,
Context
)
\rightarrow
Embodied\ Dynamical\ Modulation.\)}
\end{statementbox}

在社会情境中，它进一步表现为：

\begin{statementbox}
\centering
\adjustbox{max width=.96\linewidth}{\(\displaystyle AHC_{social}
:
(
Identity,
Authority,
InteractionHistory,
Risk,
\Psi,
Context
)
\rightarrow
EmbodiedSocialModulation.\)}
\end{statementbox}

它不保存全部社会关系，不直接产生最终行为，也不自行创造系统权限。

它承担的是一个更加基础而关键的功能：

\begin{statementbox}
将复杂的社会关系状态实时压缩成身体可以使用的低维动力学背景。
\end{statementbox}

因此，“认主”只是 AHC 具身社会绑定能力在生命周期早期的一个特殊实例。

它真正解决的问题是：

\begin{statementbox}
一个刚刚诞生的人工智能生命体，如何在尚未来得及认识世界以前，先拥有一个安全地开始认识世界的社会起点。
\end{statementbox}

而在生命周期展开以后，这个先验起点又将不断受到 $E$ 中真实经历、$\Theta$ 中长期结构预测以及 $\Psi$ 中自我关系的修正。

于是，社会关系最终不再是写入机器人数据库中的静态标签，而成为：

\begin{statementbox}
一个在多时间尺度上持续演化的具身社会动力系统。
\end{statementbox}
